\section{Conclusão}
\label{chap:conclusao}


A análise dos trabalhos discutidos evidencia a diversidade de mecanismos pelos 
quais ataques de jailbreak exploram vulnerabilidades em grandes modelos de linguagem. 
A otimização de tokens, exemplificada pelo GCG, a busca por naturalidade semântica do 
AutoDAN, as estratégias de persuasão, as transformações estruturais do FlipAttack e a
manipulação do início das respostas por prefill revelam diferentes formas de induzir 
comportamentos incompatíveis com as restrições de segurança. Em conjunto, essas 
abordagens permitem compreender a superfície de ataque como resultado da interação 
entre o funcionamento probabilístico dos modelos, suas capacidades linguísticas e os 
recursos oferecidos pelas interfaces de acesso.

Essa leitura sugere que os métodos analisados representam estratégias complementares, 
cuja combinação pode ampliar a eficácia dos ataques. A capacidade de interpretar 
instruções, reconstruir mensagens e produzir respostas coerentes, essencial à utilidade 
dos LLMs, também pode ser explorada para contornar suas salvaguardas. Nesse contexto, 
tratar o modelo como interlocutor constitui uma perspectiva analítica útil para 
investigar sua suscetibilidade a argumentos persuasivos, sem pressupor que ele 
possua processos psicológicos equivalentes aos humanos.

Do ponto de vista defensivo, os estudos discutidos apontam para a necessidade de 
avaliar conjuntamente a solicitação, o contexto e as condições de geração da resposta. 
A identificação de padrões superficiais de texto ou de pedidos explicitamente nocivos 
oferece cobertura limitada diante de ataques que preservam a fluência, ocultam a 
intenção ou manipulam a continuidade da saída. A segurança deve, portanto, considerar 
tanto o comportamento do modelo quanto os mecanismos de interação disponibilizados pelas 
aplicações.

As conclusões desta análise estão circunscritas às evidências relatadas nos 
trabalhos selecionados, sem validação experimental própria. Diferenças entre modelos, 
conjuntos de avaliação, condições de acesso e critérios de sucesso limitam comparações 
diretas e impedem estabelecer uma superioridade geral entre todas as técnicas.